\section{\texorpdfstring{The customer's view: the tables behind Sections~\ref{sec:eda} and~\ref{sec:insight}}{The customer's view: the tables behind Sections 4 and 9}}
\label{app:customer}

Everything here is computed by \code{src/customer\_view.py} from the nine tables,
the saved test predictions and the deployed lead-time model, scored on the
validation and test windows exactly as \code{src/train\_models.py} scored them.
Nothing is refitted: the rescored test forecasts equal
\code{predictions\_p1a\_test.csv} to $10^{-9}$, an assertion in the script. The
\emph{test window} is the orders placed 30 June to 30 August 2018; the
\emph{analysis window}, the orders placed 1 January 2017 to 30 August 2018. Every
figure is descriptive; the repeat-purchase comparisons are also exploratory.

\paragraph{A date the customer can rely on.} Table~\ref{tab:requote} compares
the platform's promised date with dates built from the forecast. The margin of
Section~\ref{sec:stakeholders} is chosen on the validation months alone: the
smallest margin, in steps of 0.1 day, whose broken-promise rate there is no
higher than the platform's own (4.97\%), with the date rounded up to a calendar
day --- eight days. Two margins are shown for reference: the one that matches the
platform's average quoted wait on the test window (6.9 days), chosen from quoted
dates and forecasts only, never from outcomes, and the one that matches it on the
validation months (12.4 days), which pads far more. \emph{Broken} uses the
report's definition of late (delivered after the promised date begins) for every
date; the calendar-day column counts only deliveries after the promised day has
ended.

\begin{table}[htbp]
\centering
\caption{Promised dates on the validation and test windows, from
\code{customer\_promise\_rules.csv}. Broken: share of delivered orders that arrived
after the date. Let down: broken promises among the 12{,}630 delivered test
orders. Calendar day: arrived after the promised day ended. Mean wait: the
average quoted wait, in days. Days late: mean, among broken promises. The
validation rule picks the margin; the two other margins match the platform's
average quoted wait on the test and on the validation window. ``Forecast alone''
is the forecast itself, not rounded up to a calendar date; rounded up, as every
margin is, it is broken for 25.1\% of test orders
(\code{customer\_promise\_frontier.csv}).}
\label{tab:requote}
\footnotesize
\setlength{\tabcolsep}{3.5pt}
\begin{tabular}{@{}lrrrrrr@{}}
\toprule
 & \textbf{validation} & \multicolumn{5}{c}{\textbf{test}} \\
\cmidrule(lr){2-2}\cmidrule(l){3-7}
\textbf{Date shown to the customer} & \textbf{broken} & \textbf{broken} &
\textbf{let down} & \textbf{\shortstack[r]{calendar\\day}} &
\textbf{\shortstack[r]{mean\\wait}} & \textbf{\shortstack[r]{days\\late}} \\
\midrule
Platform's own promised date & 4.97\% & 7.41\% & 936 & 4.76\% & 17.5 & 4.4 \\
Forecast alone & 26.28\% & 31.85\% & 4{,}023 & 22.37\% & 10.1 & 3.4 \\
\best{Forecast + 8 days} (validation rule) & 4.86\% & \best{2.52\%} & 318 & 2.11\% & 18.6 & 6.8 \\
Forecast + 6.9 days (platform's wait, test) & 6.03\% & 3.20\% & 404 & 2.58\% & 17.5 & 6.4 \\
Forecast + 12.4 days (platform's wait, val.) & 2.55\% & 1.10\% & 139 & 0.92\% & 23.0 & 8.5 \\
\bottomrule
\end{tabular}
\end{table}

By region the margin moves the padding to where the risk is
(\code{customer\_\allowbreak promise\_\allowbreak region.csv}). Southeast
customers (8{,}951 test orders)
are let down 8.59\% of the time by the platform's date and 2.41\% by the forecast
plus eight days, for an average wait quoted 1.4 days longer; in the North and
Northeast (1{,}312) the rate falls from 5.03\% to 3.58\% at the same average wait
(26.6 days); in the South and Central-West (2{,}367) from 4.27\% to 2.32\%, for
0.6 days more. Overall 61.5\% of customers would see a later date than today and
31.3\% an earlier one, by at least a day, and the dates still missed are missed
by more: 6.8 days on average, against 4.4. Three cautions. The validation months
include the May 2018 strike and June's padded promises, so the margin
over-delivers on the test window: 2.52\% broken, against 7.41\% for the
platform's date there and the 4.97\% bar it was chosen to meet on validation.
The forecast over-predicts by 1.74 days on the test window and the margin
absorbs that bias, so it must be re-chosen at every refit. And the date assumes
that delivery times and customers do not change because the quote does; the
forecast also takes the platform's own estimate as an input, so the platform
would keep computing that estimate and change only the date it shows.

\paragraph{What happened to every 1{,}000 orders.} Table~\ref{tab:outcomes} is
the full count behind Table~\ref{tab:customer}, with 95\% Wilson intervals
(\code{customer\_outcomes\_test.csv}). Every row names its population: orders
placed in the test window (12{,}927) or, for the checkout warning, delivered test
orders (12{,}630).

\begin{table}[htbp]
\centering
\caption{What happened to the customers of the test window, from
\code{customer\_outcomes\_test.csv}. Per 1{,}000 orders of the population named
in the third column, with 95\% Wilson intervals. Indented rows are subsets of
the row above them. ``Reachable'' names the contact that could still have come
before the customer rated the order (Appendix~\ref{app:timing}).}
\label{tab:outcomes}
\footnotesize
\setlength{\tabcolsep}{4pt}
\begin{tabular}{@{}p{8.2cm}rlrl@{}}
\toprule
\textbf{Customers whose order\ldots} & \textbf{orders} & \textbf{of} & \textbf{per 1{,}000} & \textbf{95\% interval} \\
\midrule
was placed (purchases 30 June to 30 August 2018) & 12{,}927 & placed & 1{,}000.0 & --- \\
parcel arrived by the promised date & 11{,}694 & placed & 904.6 & [899.4, 909.6] \\
parcel arrived after the promised date began (the report's definition) & 936 & placed & 72.4 & [68.1, 77.0] \\
\quad during the promised day itself & 335 & placed & 25.9 & [23.3, 28.8] \\
\quad after the promised day ended & 601 & placed & 46.5 & [43.0, 50.3] \\
parcel never delivered within the extract & 297 & placed & 23.0 & [20.5, 25.7] \\
\quad order cancelled or unavailable & 149 & placed & 11.5 & [9.8, 13.5] \\
\addlinespace[2pt]
no parcel when the promised day ended & 898 & placed & 69.5 & [65.2, 74.0] \\
\quad of whom never received it & 297 & placed & 23.0 & [20.5, 25.7] \\
\addlinespace[2pt]
rated the order 1--2 stars & 1{,}418 & placed & 109.7 & [104.4, 115.2] \\
\quad parcel on time & 923 & placed & 71.4 & [67.1, 76.0] \\
\quad parcel late & 287 & placed & 22.2 & [19.8, 24.9] \\
\quad parcel never delivered & 208 & placed & 16.1 & [14.1, 18.4] \\
late parcel, 1--2 stars written before it arrived & 222 & placed & 17.2 & [15.1, 19.6] \\
no parcel when the promised day ended, not yet rated then, went on to rate 1--2 stars & 440 & placed & 34.0 & [31.0, 37.3] \\
\addlinespace[2pt]
rated 1--2 stars; reachable only by the promised-date message & 398 & placed & 30.8 & [27.9, 33.9] \\
rated 1--2 stars; reachable only by a call at delivery & 945 & placed & 73.1 & [68.7, 77.7] \\
rated 1--2 stars; reachable by either contact & 42 & placed & 3.2 & [2.4, 4.4] \\
rated 1--2 stars before either moment & 33 & placed & 2.6 & [1.8, 3.6] \\
\addlinespace[2pt]
on the slides' daily top-ten list (grouped by purchase day) & 614 & placed & 47.5 & [44.0, 51.3] \\
\quad and rated the order 1--2 stars & 351 & placed & 27.2 & [24.5, 30.1] \\
\quad of whom not yet rated when the promised day ended (no parcel then) & 271 & placed & 21.0 & [18.6, 23.6] \\
\quad of whom reachable only by a call at delivery & 48 & placed & 3.7 & [2.8, 4.9] \\
\addlinespace[2pt]
flagged at checkout (late risk at or above the 11\% line) & 5{,}647 & delivered & 447.1 & [438.5, 455.8] \\
\quad parcel then late & 773 & delivered & 61.2 & [57.2, 65.5] \\
late parcel, not flagged & 163 & delivered & 12.9 & [11.1, 15.0] \\
flagged, parcel then on time & 4{,}874 & delivered & 385.9 & [377.5, 394.4] \\
\bottomrule
\end{tabular}
\end{table}

A parcel delivered during the promised day counts as late under the report's
definition: 335 of the 936 late test parcels. Their customers rate them much like
on-time ones (Table~\ref{tab:lateday}, analysis window): 12.3\% gave 1--2 stars,
against 11.1\% for delivery during the day before and 19.8\% for delivery a day
later. That is why the promised-date message is triggered when the promised day
ends, not when it begins.

\begin{table}[htbp]
\centering
\caption{The 1--2 star rate by how early or late the parcel arrived against the
promised date, on the delivered and reviewed orders of the analysis window, with
95\% Wilson intervals (\code{customer\_lateness\_days.csv}). The promised date is
a calendar date, so ``during the promised day'' is up to one day after it began.}
\label{tab:lateday}
\footnotesize
\begin{tabular}{@{}lrrl@{}}
\toprule
\textbf{Parcel arrived} & \textbf{orders} & \textbf{1--2 stars} & \textbf{95\% interval} \\
\midrule
7 or more days early & 70{,}673 & 8.9\% & [8.7, 9.2] \\
3--7 days early & 12{,}527 & 9.8\% & [9.3, 10.4] \\
1--3 days early & 3{,}253 & 11.3\% & [10.2, 12.4] \\
up to 1 day early & 1{,}449 & 11.1\% & [9.6, 12.8] \\
during the promised day & 1{,}280 & 12.3\% & [10.7, 14.3] \\
1--2 days late & 819 & 19.8\% & [17.2, 22.6] \\
2--3 days late & 536 & 34.3\% & [30.4, 38.4] \\
3--7 days late & 1{,}773 & 61.3\% & [59.0, 63.5] \\
more than 7 days late & 3{,}251 & 78.3\% & [76.8, 79.7] \\
\bottomrule
\end{tabular}
\end{table}

\paragraph{What unhappy customers wrote.} The comments of the 1--2 star reviews
of the analysis window (the review the repository keeps per order, title and
message together, lower-cased, accents stripped and common abbreviations
expanded) are matched against seven overlapping topics, each a regular
expression in \code{src/customer\_view.py}, listed in
\code{customer\_complaint\_dictionary.csv}. The patterns were frozen before their
precision was measured on a random sample of 25 matched comments per topic, each
comment read and marked right or wrong by the AI assistant used for the revision
(AI Tool Declaration); \code{customer\_topic\_labels.csv} records the 175
verdicts, with a note on each miss, so any of them can be re-read. Table~\ref{tab:topics} gives the
shares. They are lower bounds: 18.9\% of comments match no topic, often a
paraphrase the patterns miss. The comments separate two failures. Customers whose
parcel was late or never came write that they do not have it (59.9\% and 58.3\%),
and those whose order never came also write about refunds and cancellation
(24.5\%); customers whose parcel came on time write about what was in it: a wrong
or missing item (34.0\%), quality or a product unlike its description (16.9\%),
damage or a defect (15.0\%). Comments are written with the rating, so they are
read here and never used as a feature.

\begin{table}[htbp]
\centering
\caption{Topics of 1--2 star comments on the analysis window, as \% of the
comments in each group (a comment can match several topics), and the held-out
precision of each pattern with its 95\% Wilson interval. From
\code{customer\_complaint\_topics.csv} and \code{customer\_complaint\_precision.csv}.}
\label{tab:topics}
\footnotesize
\setlength{\tabcolsep}{4pt}
\begin{tabular}{@{}lrrrrr@{}}
\toprule
 & \multicolumn{4}{c}{\textbf{the parcel was}} & \\
\cmidrule(lr){2-5}
\textbf{Topic} & \textbf{any} & \textbf{on time} & \textbf{late} & \textbf{never delivered} & \textbf{precision} \\
\midrule
1--2 star reviews & 14{,}365 & 8{,}084 & 4{,}136 & 2{,}145 & \\
\quad with a comment & 10{,}826 & 6{,}385 & 2{,}889 & 1{,}552 & \\
\midrule
not received (yet) & 34.5 & 17.3 & 59.9 & 58.3 & 88 [70, 96] \\
late, slow, missed date & 12.4 & 5.4 & 24.6 & 18.2 & 96 [80, 99] \\
wrong or missing item & 22.2 & 34.0 & 5.5 & 4.3 & 88 [70, 96] \\
defective or damaged & 9.5 & 15.0 & 1.7 & 1.8 & 92 [75, 98] \\
quality, not as described & 10.6 & 16.9 & 2.1 & 0.5 & 96 [80, 99] \\
seller service, communication & 13.9 & 12.0 & 16.3 & 17.5 & 80 [61, 91] \\
refund, cancel, return & 14.2 & 14.8 & 7.4 & 24.5 & 96 [80, 99] \\
no topic matched & 18.9 & 20.2 & 17.0 & 17.5 & --- \\
\bottomrule
\end{tabular}
\end{table}

\paragraph{Orders with several items or sellers.} On the 87{,}902 on-time,
delivered and reviewed orders of the analysis window (Table~\ref{tab:basket}), a
single item from one seller draws 1--2 stars 7.5\% of the time, two or more items
from one seller 21.0\%, and two or more sellers 47.1\%. Multi-item orders are
10.1\% of on-time orders but 27.1\% of their 1--2 star reviews, and 74.9\% and
84.2\% of those reviewers who comment say that part of the order was wrong,
missing or never came. Two limits: \code{n\_items} counts item rows, so two units
of one product count as two; and an order carries one delivery timestamp, so
``on time'' does not mean that every package arrived. The finding is an
association, not a mechanism.

\begin{table}[htbp]
\centering
\caption{On-time, delivered and reviewed orders of the analysis window by
basket, from \code{customer\_basket.csv}: the 1--2 star rate with its 95\% Wilson
interval, the basket's share of all on-time 1--2 star reviews, and what the
reviewers who commented wrote (the topics of Table~\ref{tab:topics}; the first
combines ``wrong or missing item'' and ``not received'').}
\label{tab:basket}
\footnotesize
\setlength{\tabcolsep}{4pt}
\begin{tabular}{@{}p{6.6cm}rrrr@{}}
\toprule
 & \textbf{1 item,} & \textbf{2+ items,} & & \textbf{all} \\
 & \textbf{1 seller} & \textbf{1 seller} & \textbf{2+ sellers} & \textbf{on time} \\
\midrule
Orders & 79{,}011 & 7{,}650 & 1{,}241 & 87{,}902 \\
\quad share of on-time orders & 89.9\% & 8.7\% & 1.4\% & 100.0\% \\
Rated 1--2 stars & 7.5\% & 21.0\% & 47.1\% & 9.2\% \\
\quad 95\% interval & [7.3, 7.6] & [20.1, 21.9] & [44.4, 49.9] & [9.0, 9.4] \\
\quad share of all on-time 1--2 star reviews & 72.9\% & 19.9\% & 7.2\% & 100.0\% \\
1--2 star reviews with a comment & 4{,}482 & 1{,}403 & 500 & 6{,}385 \\
\quad say part of the order was wrong, missing or never came & 33.0\% & 74.9\% & 84.2\% & 46.2\% \\
\quad say it was defective or damaged & 19.3\% & 5.5\% & 2.4\% & 15.0\% \\
\quad say its quality was poor or not as described & 21.7\% & 7.0\% & 2.4\% & 16.9\% \\
\bottomrule
\end{tabular}
\end{table}

\paragraph{Do let-down customers come back less?} Each customer
(\code{customer\_\allowbreak unique\_\allowbreak id}) is followed from their first
order. A repeat is another order
placed after the first was delivered and within the horizon of that delivery;
the window must close by 22 August 2018, after which the extract thins out
(\code{customer\_repeat.csv}, Table~\ref{tab:repeat}). Of the 50{,}360 first-time
buyers that leaves at 180 days, 1.84\% came back. After a late first parcel
1.44\% did, against 1.86\% after an on-time one; after a 1--2 star first review
1.50\%, against 1.90\% after 4--5 stars. The direction is the same at every
horizon, and some differences are nominally significant (after a 1--2 star review
at 90 days, $-0.36$ points, $p = 0.003$), but they are tenths of a point on
repeat rates of 0.56\% to 3.21\%, close to the smallest difference this sample
can detect, and 28 comparisons were run over four horizons and two definitions
of a repeat.
They are exploratory. A bad experience therefore cannot be priced through the
customer's own return, which is why the R\$40 of Section~\ref{sec:cost} is an
assumption standing for the rating's effect on later buyers.

\begin{table}[htbp]
\centering
\caption{Repeat purchase by the first order's experience, from
\code{customer\_repeat.csv}: first-time buyers whose first order was delivered
at least the horizon before 22 August 2018; a repeat is an order placed after
that delivery and within the horizon. Rates in \% (``all'': every buyer of the
cohort); the difference is in percentage points with a 95\% Newcombe interval;
the odds ratio is Mantel--Haenszel, stratified by customer region and purchase
quarter, with its 95\% interval. The other definition (any later order on a later
day, counted from purchase) is in the CSV.}
\label{tab:repeat}
\footnotesize
\setlength{\tabcolsep}{4pt}
\begin{tabular}{@{}lrrrrll@{}}
\toprule
 & \textbf{days} & \textbf{buyers} & \textbf{all} & \textbf{group vs reference} &
\textbf{difference} & \textbf{odds ratio} \\
\midrule
\multicolumn{7}{@{}l}{\textit{First parcel late, against on time}} \\
 & 30 & 82{,}422 & 0.56 & 0.41 vs 0.57 & $-$0.16 [$-$0.30, $+$0.03] & 0.79 [0.53, 1.18] \\
 & 90 & 72{,}007 & 1.16 & 0.86 vs 1.19 & $-$0.33 [$-$0.55, $-$0.05] & 0.79 [0.59, 1.06] \\
 & 180 & 50{,}360 & 1.84 & 1.44 vs 1.86 & $-$0.43 [$-$0.82, $+$0.08] & 0.83 [0.61, 1.13] \\
 & 365 & 18{,}546 & 3.21 & 2.06 vs 3.26 & $-$1.20 [$-$2.07, $+$0.20] & 0.64 [0.37, 1.09] \\
\multicolumn{7}{@{}l}{\textit{First review 1--2 stars, against 4--5 stars}} \\
 & 30 & 82{,}422 & 0.56 & 0.42 vs 0.59 & $-$0.17 [$-$0.29, $-$0.02] & 0.75 [0.55, 1.02] \\
 & 90 & 72{,}007 & 1.16 & 0.88 vs 1.23 & $-$0.36 [$-$0.55, $-$0.13] & 0.74 [0.59, 0.93] \\
 & 180 & 50{,}360 & 1.84 & 1.50 vs 1.90 & $-$0.40 [$-$0.71, $-$0.03] & 0.81 [0.65, 1.01] \\
 & 365 & 18{,}546 & 3.21 & 2.87 vs 3.33 & $-$0.46 [$-$1.18, $+$0.43] & 0.86 [0.65, 1.14] \\
\bottomrule
\end{tabular}
\end{table}

\paragraph{The value of a contact under other assumptions.} Table~\ref{tab:value}
restates each list in reais per 1{,}000 reviewed test orders (12{,}855) at the
assumed R\$5 a contact and R\$40 a 1--2 star rating, for three shares of the
reached customers a contact turns around: all, half and a quarter. Nothing in
the data measures that share. \code{customer\_value\_sensitivity.csv} also varies
the contact cost (R\$2.5, 5, 10) and the rating's value (R\$20, 40, 80); only
the share times that value counts, so the ``half'' column is also every reached
customer turned around at R\$20 a rating. Read the other way, at R\$5 a contact
the promised-date message pays at full effect if a 1--2 star rating is worth at
least R\$9.1, the delivery-day call R\$19.2 (R\$36.5 and R\$76.9 at one in four;
\code{customer\_value.csv}). The
first panel counts every hit, as the slides did, including customers who had
already rated the order or whose order never came; the second counts only the
customers each contact could still reach at its own moment. The two contacts
Section~\ref{sec:stakeholders} recommends are the last two rows.

\begin{table}[htbp]
\centering
\caption{Net value per 1{,}000 reviewed test orders at the assumed R\$5 : R\$40,
by the share of reached 1--2 star customers a contact turns around, from
\code{customer\_value.csv}. Hits: listed customers who went on to rate the order
1--2 stars. ``Must turn around'': the break-even share; above 100\% the contact
cannot pay even if it turns every reached customer around.}
\label{tab:value}
\footnotesize
\setlength{\tabcolsep}{4pt}
\begin{tabular}{@{}lrrrrrrr@{}}
\toprule
 & & & & \multicolumn{3}{c}{\textbf{share turned around}} & \textbf{must turn} \\
\cmidrule(lr){5-7}
\textbf{List} & \textbf{listed} & \textbf{hits} & \textbf{precision} & \textbf{all} & \textbf{half} & \textbf{a quarter} & \textbf{around} \\
\midrule
\multicolumn{8}{@{}l}{\textit{Every hit counted, as on the slides}} \\
Slides: threshold 0.11 & 1{,}927 & 712 & 0.369 & 1{,}466 & 358 & $-$196 & 33.8\% \\
Slides: ten a day by purchase day & 614 & 351 & 0.572 & 853 & 307 & 34 & 21.9\% \\
Deployed threshold 0.19 & 981 & 564 & 0.575 & 1{,}373 & 496 & 57 & 21.7\% \\
\midrule
\multicolumn{8}{@{}l}{\textit{Only customers the contact can still reach at its moment}} \\
Threshold 0.11, call at delivery & 1{,}927 & 282 & 0.146 & 128 & $-$311 & $-$530 & 85.4\% \\
Slides' ten a day, call at delivery & 614 & 54 & 0.088 & $-$71 & $-$155 & $-$197 & 142.1\% \\
Threshold 0.19, call at delivery & 981 & 137 & 0.140 & 45 & $-$168 & $-$275 & 89.5\% \\
Promised-date message & 804 & 440 & 0.547 & 1{,}056 & 372 & 30 & 22.8\% \\
Delivery-day call, ten a day & 546 & 142 & 0.260 & 229 & 9 & $-$102 & 48.1\% \\
\bottomrule
\end{tabular}
\end{table}

Table~\ref{tab:decomp} splits the figures quoted on the slides by whom a contact
would have reached. Of the slides' threshold-0.11 list, R\$877 per 1{,}000 of
gross value comes from delivered customers who had not yet written, R\$647 from
orders that never arrived and R\$691 from late parcels already rated while the
customer waited, against R\$750 of contacts: 60.4\% of the gross value is
customers whose order never came or who rated a late parcel before it arrived. A call at delivery alone nets
R\$128 on that list and $-$R\$71 on the ten-a-day list.

\begin{table}[htbp]
\centering
\caption{Where the slides' values come from, per 1{,}000 reviewed test orders at
the assumed R\$5 : R\$40 with every contact counted as turning its customer
around, from \code{customer\_value\_decomposition.csv}. ``Not written'': the hits
who had not rated the order before the parcel arrived, plus every hit whose
order never arrived, whenever they rated (the slides' R\$775 and R\$433).}
\label{tab:decomp}
\footnotesize
\begin{tabular}{@{}lrrr@{}}
\toprule
 & \textbf{threshold 0.11} & \textbf{threshold 0.19} & \textbf{ten a day} \\
 & \textbf{(slides)} & \textbf{(deployed)} & \textbf{(slides)} \\
\midrule
Customers listed & 1{,}927 & 981 & 614 \\
Hits: delivered, not yet written & 282 & 137 & 54 \\
\quad never delivered & 208 & 208 & 162 \\
\quad late, written before arrival & 222 & 219 & 135 \\
\midrule
Gross R\$: delivered, not yet written & 877 & 426 & 168 \\
\quad never delivered & 647 & 647 & 504 \\
\quad late, written before arrival & 691 & 681 & 420 \\
Contacts at R\$5 & $-$750 & $-$382 & $-$239 \\
\midrule
Net, every hit counted & 1{,}466 & 1{,}373 & 853 \\
Net, not written before arrival & 775 & 692 & 433 \\
Net, a call at delivery alone & 128 & 45 & $-$71 \\
Gross share: undelivered or written first & 60.4\% & 75.7\% & 84.6\% \\
\bottomrule
\end{tabular}
\end{table}
